\section{Substitution and strong normalization for simple types}
\label{sec:p07-stlc-metatheory}

The arrow-only calculus on Sheet~\ref{sheet:stlc} has a termination guarantee
for every choice of beta redex. Its proof uses a family of predicates indexed
by types, rather than a count of term symbols: substitution can duplicate an
argument. The following treatment supplies the substitution and reducibility
arguments for the calculus in Section~\ref{sec:normalization-curry-howard-02}.
The reducibility construction follows the arrow fragment of
\cite[Chapter~6]{GirardLafontTaylor2003Web}.

\begin{hbdefinition}{def:p07-stlc-meta-setting}{Terms, contexts, and substitutions}
Types are generated from a set of base types by
\(\alpha\mathbin{\to_{\mathrm{ty}}}\beta\). Terms are finite variables,
applications \(MN\), and abstractions \(\lambda x{:}\alpha.M\), taken
modulo consistent renaming of bound variables. A typing context \(\Theta\)
is a finite list assigning one type to each of its distinct variables.
Write \(\Theta\vdash_{\mathsf{STLC}}M:\alpha\) for exactly the three
typing rules on Sheet~\ref{sheet:stlc}. There are arbitrarily many fresh
variable names. A context may be enlarged by assignments to fresh variables
or reordered without changing a derivable judgment.

A simultaneous substitution \(\rho\) assigns a term to each variable of
a finite domain and acts as the identity on other free variables. Its action
\(M\rho\) traverses the original term once; it does not recursively
substitute into the inserted images. Before crossing an abstraction, rename
its binder fresh for the domain and the free variables of all substitution
images, and retain its type annotation. In particular, \(M[N/x]\) replaces
the free occurrences of \(x\), avoiding capture.

Write \(\to_\beta\) for a beta contraction at any term position, including
inside an abstraction, and \(\to_\beta^*\) for its finite reflexive
transitive closure. A term is \emph{strongly normalizing} if no infinite
\(\to_\beta\)-sequence starts from it. A beta-normal term has no beta
redex. A \emph{maximal} reduction sequence is either infinite or ends at a
term with no successor. This last qualification excludes stopping a finite
sequence early at a reducible term.
\index{simply typed lambda calculus!substitution}
\index{strong normalization!simple types}
\end{hbdefinition}

\begin{hblemma}{thm:p07-stlc-substitution}
 {proved=annex:p7.stlc.substitution}{Simultaneous typing substitution}
Suppose \(\Theta\vdash_{\mathsf{STLC}}M:\alpha\), and for every
\(x:\beta\) in \(\Theta\) a common target context \(\Delta\) derives
\(\Delta\vdash_{\mathsf{STLC}}\rho(x):\beta\). Then
\(\Delta\vdash_{\mathsf{STLC}}M\rho:\alpha\).
In particular, if \(\Theta,x:\beta\vdash_{\mathsf{STLC}}M:\alpha\)
and \(\Theta\vdash_{\mathsf{STLC}}N:\beta\), then
\(\Theta\vdash_{\mathsf{STLC}}M[N/x]:\alpha\).
\end{hblemma}

\begin{hbtheorem}{thm:p07-stlc-subject-reduction}
 {proved=annex:p7.stlc.subject-reduction}{Subject reduction for simple types}
If \(\Theta\vdash_{\mathsf{STLC}}M:\alpha\) and
\(M\to_\beta N\), then \(\Theta\vdash_{\mathsf{STLC}}N:\alpha\).
Consequently, every finite beta reduct has the same type in the same context.
\end{hbtheorem}

\begin{hbtheorem}{thm:p07-stlc-strong-normalization}
 {proved=annex:p7.stlc.strong-normalization}{Strong normalization for simple types}
Every term typable on Sheet~\ref{sheet:stlc}, including an open term in a
finite typing context, is strongly normalizing for compatible beta reduction.
Every maximal beta-reduction sequence from it is therefore finite and ends
at a beta-normal term of the same type.
\end{hbtheorem}

\begin{hbdefinition}{def:p07-stlc-reducibility}{Reducibility predicates}
For the proof, consider all raw finite terms with the same annotated syntax,
without requiring them to have a typing derivation. Let
\(\mathsf{SN}_\beta\) be the set of strongly normalizing raw terms.
Define sets \(\mathcal R_\alpha\) by recursion on the type:
\[
 \mathcal R_{\mathsf b}=\mathsf{SN}_\beta,\qquad
 \mathcal R_{\alpha\mathbin{\to_{\mathrm{ty}}}\beta}
 =\{M:\text{ for every }N\in\mathcal R_\alpha,
                       \ MN\in\mathcal R_\beta\}.
\]
These auxiliary sets may contain untypable terms. The fundamental lemma in
the annex places a term in the predicate indexed by its actual derived type.
A raw term is \emph{neutral} here when it is not an abstraction. In particular,
an application with an abstraction at its head is neutral, even though it
has a beta redex. Neutrality and normality are different properties.
\index{reducibility predicate}\index{neutral term!normalization proof}
\end{hbdefinition}

\begin{hbcorollary}{thm:p07-stlc-inert-constants}
 {proved=annex:p7.stlc.inert-constants}{Adding typed inert constants}
Extend the arrow-only term syntax by constants, each with one specified simple
type, and a typing rule admitting a constant of that type in every context.
If constants have no reduction rules, every well-typed term of the extension
is strongly normalizing for compatible beta reduction. No finiteness of the
whole constant signature is required.
\end{hbcorollary}

\begin{hbmethod}{def:p07-stlc-normalizer}{Computing a normal form}
For a term with a typing derivation, repeatedly traverse its finite syntax
tree in a fixed order and contract the first beta redex found. Capture-avoiding
substitution is the operation in Definition~\ref{def:p07-stlc-meta-setting}.
Stop when the traversal finds none. Theorem~\ref{thm:p07-stlc-strong-normalization}
proves that the procedure stops; Theorem~\ref{thm:p07-stlc-subject-reduction}
preserves its type. The same termination conclusion holds for any strategy
that contracts a beta redex whenever one exists.

Termination alone does not establish independence of the chosen strategy.
For that conclusion use the Church--Rosser theorem
\ref{thm:p07-church-rosser}: two reached normal forms must be
alpha-equivalent. Normalizing both sides therefore decides beta equivalence
of well-typed inputs, because beta-convertible terms have a common reduct,
and normal terms admit only the reflexive reduction. This algorithmic
conclusion assumes finite input terms and effective syntax operations; it
does not assert a small running-time bound.
\end{hbmethod}

\begin{hbexample}{def:p07-stlc-normalization-example}{Argument duplication and a type boundary}
Let \(x:\alpha\), \(f:\alpha\mathbin{\to_{\mathrm{ty}}}\alpha\), and
\(I=\lambda z{:}\alpha.z\). The term
\[
 (\lambda g{:}\alpha\mathbin{\to_{\mathrm{ty}}}\alpha.\,g(gx))
       (\lambda y{:}\alpha.\,f(Iy))
\]
has type \(\alpha\). Contracting the outer redex duplicates its argument
and gives
\[
 (\lambda y{:}\alpha.\,f(Iy))
       ((\lambda y{:}\alpha.\,f(Iy))x)
 \to_\beta^* f(fx).
\]
Alternatively, first reducing the occurrence of \(Iy\) inside the argument
gives \((\lambda g.\,g(gx))(\lambda y.\,fy)\), with the same annotations
as above, and then \(f(fx)\). This example explains why the proof permits
reductions under binders and does not use term size as a decreasing measure.

For comparison, the raw untyped term
\(\Omega=(\lambda x.xx)(\lambda x.xx)\) reduces to itself.
Self-application would require one variable type both to be \(\delta\)
and \(\delta\mathbin{\to_{\mathrm{ty}}}\gamma\). No finite simple type
satisfies that equation: counting type constructors makes the right side
strictly larger. Thus this divergence supplies no well-typed counterexample
to the theorem. Adding an operator with a new unfolding rule requires a new
normalization argument; the inert-constant corollary concerns constants with
no such rule.
\end{hbexample}
